\section{Quantifying and Predicting Legal Outcomes (2014--2018)}
\label{sec:outcomes}

\subsection{2014: predicting a court's decision from data}

The next advance did not come from representing legal reasoning more carefully. It came from a
different question: could statistics predict a court's decision, with no explicit model of legal
reasoning at all? In 2014, \citet{katz2014} built a time-evolving statistical classifier, training it
on historical voting and case data to predict the outcomes of United States Supreme Court decisions.
Their method used an ensemble of decision trees, a random-forest classifier; it was not a language
model, and it was not a case-based comparison in Ashley's sense, but rather a quantitative
pattern-recognition problem over structured case features. The result mattered precisely because the
system did not need to represent legal rules or compare case facts the way HYPO did: it showed that
a purely statistical model, given enough historical data, could forecast judicial behavior with
useful accuracy. This also exposed a limitation that recurs throughout the rest of this survey. A
system can predict an outcome correctly while giving no account of the legal reasoning behind it,
which is a very different capability from explaining \emph{why} a court should decide a case a
particular way.

In the same year, a separate study took on a different question: not whether a system could predict
an outcome, but whether a system could answer a legal question directly. \citet{kim2014} built a
legal question-answering system that combines legal information retrieval with textual entailment,
ranking candidate answers with a support vector machine informed by syntactic and semantic
similarity. The authors evaluated it on the first Competition on Legal Information Extraction and
Entailment (COLIEE) dataset. Katz predicted an outcome from structured case data, while Kim retrieved
and compared legal text directly, and together they give a useful snapshot of 2014: the same year
produced an outcome-prediction system that bypassed legal text entirely, and a question-answering
system built entirely around retrieving and comparing it. The COLIEE benchmark this study helped
establish recurs throughout this survey, as later editions appear in subsequent years.

\subsection{2017: extracting structure from legal text, and predicting charges from facts}

Three years later, two separate lines of work, both published in 2017, pushed in a complementary
direction: instead of predicting an outcome from case metadata, they tried to connect a predicted
outcome back to the specific legal text that justifies it. \citet{chalkidis2017} took on contract
element extraction, applying neural sequence-labelling methods to identify defined terms,
obligations, and other structured elements directly within contract text. This was an early sign
that neural NLP methods could extract the kind of structured legal information earlier rule-based
systems had encoded by hand. In the same year, \citet{luo2017} proposed a model for predicting
criminal charges that explicitly retrieves and attends over the specific statutory articles that
justify a predicted charge. It uses an attention mechanism to connect the facts of a case to the law
articles it identifies as relevant, and the authors report that this explicit legal-basis signal
improved charge-prediction accuracy over models using case facts alone. Katz's 2014 model predicted
an outcome without reference to legal text, while Luo et al.'s model took a first step toward
grounding a predicted outcome in the specific statutory provisions that justify it. This anticipates
the evidence-grounding problem that motivates legal RAG later in this survey.

\subsection{2018: modelling legal judgment as several connected subtasks}

The following year, \citet{zhong2018} generalized this idea. Rather than predicting a single output,
a charge or a verdict, they observed that a real legal judgment involves several connected subtasks:
identifying the applicable law articles, determining the charge, and deciding the sentence term.
They modelled the dependencies among these subtasks as a directed acyclic graph and solved them
jointly through topological multi-task learning, reporting that representing these dependencies
explicitly improved performance over treating each subtask on its own. Compared with
\citet{katz2014} and \citet{luo2017}, this period teaches one specific and cumulative lesson: legal
judgment prediction benefits from representing the structure of a legal decision, not just its
textual facts. The rest of this survey shows the field relearning that lesson, in a different form,
as it moves from outcome prediction toward retrieval-augmented generation.

\subsection{What the foundational and prediction periods together established}

Table~\ref{tab:foundational} summarizes the approaches discussed so far, in the order the field
published them between 1991 and 2018.

\begin{table*}[tp]
\centering
\caption{Legal AI milestones discussed so far, in chronological order (1991--2018).}
\label{tab:foundational}
\footnotesize
\setlength{\tabcolsep}{5pt}
\renewcommand{\arraystretch}{1.25}
\begin{tabularx}{\textwidth}{>{\raggedright\arraybackslash}p{1.6cm} >{\raggedright\arraybackslash}p{2.9cm} >{\raggedright\arraybackslash}X >{\raggedright\arraybackslash}X}
\toprule
\textbf{Year} & \textbf{Study} & \textbf{Core idea} & \textbf{Principal limitation} \\
\midrule
1991 & \citet{ashley1991} (HYPO) & Compare a dispute with precedents and construct arguments from relevant similarities and differences. & Legally meaningful similarity, not textual similarity, must be identified. \\
2003 & \citet{ashley2003} & Reframe case-based legal reasoning as lazy, case-driven learning. & Does not by itself resolve how legal knowledge should be found or maintained at scale. \\
2009 & \citet{cormack2009} & Combine multiple ranked lists via Reciprocal Rank Fusion (general-purpose IR, not legal). & Carries no legal content; legal relevance only emerges once applied to legal retrieval, later in this survey. \\
2014 & \citet{katz2014} & Predict United States Supreme Court outcomes with a statistical, time-evolving classifier. & Predicts outcomes without explaining the applicable legal reasoning. \\
2014 & \citet{kim2014} & Answer legal questions by combining information retrieval with SVM-ranked textual entailment on COLIEE data. & Entailment between a question and a retrieved article, not broader legal reasoning. \\
2017 & \citet{chalkidis2017} & Extract structured elements (terms, obligations) from contracts with neural sequence labelling. & Extraction, not reasoning about the extracted elements. \\
2017 & \citet{luo2017} & Jointly predict criminal charges and the statutory articles that justify them. & Evaluated on Chinese criminal cases; generalization across jurisdictions untested. \\
2018 & \citet{zhong2018} & Model law article, charge, and sentence prediction as dependent subtasks solved jointly. & Predicts a judgment's components without retrieving or citing external authoritative sources. \\
\bottomrule
\end{tabularx}
\end{table*}

None of these approaches alone solved the problem of locating authoritative evidence, keeping its
context, and producing a legally justified answer. Together, though, they set up the two remaining
halves of the problem, and the rest of this survey follows them in order. The first half is
representing legal knowledge in a form a model can learn from, which the field pursued through
legal language models starting in 2020. The second half is retrieving external, verifiable evidence
instead of relying on what a model has memorized, which the field pursued through
retrieval-augmented generation starting in 2024. We take up the first of these next.
